\chapter{Research Methodology}
\label{ch:methodology}

\section{Introduction}
The project combines experimental retrieval comparisons with iterative software
development. The research question is whether communicating the intended garment
before image encoding improves retrieval under a defined relevance criterion.
Supporting questions concern the choice of pretrained encoder, the effect of
category filtering, and whether modifying gallery images justifies replacing the
original index. The engineering question is whether the resulting workflow can
be used and maintained through a local browser application.

These questions require different evidence. Labelled retrieval queries measure
ranking behaviour; unit and integration tests exercise software contracts;
browser observations examine visible interaction; reviewed alternate views
provide limited identity evidence. The study records these evidence types
separately rather than combining them into a single accuracy figure.

\section{Data Collection and Dataset Specification}
\subsection{Source and Partition Roles}
The image collection is the public ten-class \texttt{clothing-dataset-small}
subset~\cite{GrigorevClothingDataset}. Its local inventory contains 3,781 images:
3,068 in the named training split, 341 in validation, and 372 in test.
Table~\ref{tab:data-splits} gives the category distribution. The original directory
key \texttt{outwear} is retained in code for compatibility and described as
outerwear in the report.

\begin{table}[H]
\centering\caption{Verified image inventory by category and partition}
\label{tab:data-splits}
\begin{ReportTable}% Generated by generate_evidence.py; edit the source data or script.
\begin{tabular}{@{}lrrr@{}}
\toprule
\textbf{Category} & \textbf{Gallery} & \textbf{Validation} & \textbf{Test} \\ \midrule
Dress & 241 & 32 & 15 \\
Hat & 123 & 14 & 12 \\
Long Sleeve & 455 & 49 & 72 \\
Outerwear & 184 & 24 & 38 \\
Pants & 468 & 49 & 42 \\
Shirt & 290 & 29 & 26 \\
Shoes & 198 & 26 & 73 \\
Shorts & 202 & 25 & 30 \\
Skirt & 112 & 12 & 12 \\
T-Shirt & 795 & 81 & 52 \\
Total & 3068 & 341 & 372 \\
\bottomrule
\end{tabular}
\end{ReportTable}
\end{table}

The training split is the retrieval gallery, not a newly collected training
set for the neural models. All encoders used in this project are pretrained.
Validation supports comparison and development choices; test supplies the
retained final category comparisons. Validation and test images are not inserted
into the application's original research gallery.

\subsection{Identity and Duplicate Auditing}
The supplied dataset does not include verified merchant SKU identities. The
application derives an image identifier from each filename and maps it to a
category and display title. Two different filenames may still depict the same
garment, and two similar garments may be different products. Consequently,
3,068 is a count of searchable images, not a verified count of unique products.

The recorded dataset audit computed both file hashes and decoded, EXIF-aware
RGB-pixel hashes across all 3,781 images. It found no exact byte or decoded-pixel
duplicate groups. A fresh byte-level audit during report preparation again found
no overlap between the three partitions. These checks do not rule out resized,
recompressed, cropped, or alternate-view photographs of the same garment.
Visual near-duplicate judgments are therefore kept distinct from exact hash
matches and from independently verified commercial identities.

\subsection{Independent Added-Product Data}
Products entered through the management page form a second catalogue. Each has
an explicit product UUID and one to five associated photographs. Their vectors
are produced by the same deployed encoder but stored with product data in a
separate SQLite database. This catalogue supports application use and CRUD
testing; it is not added to the research gallery or silently included in the
historical evaluation results.

\section{Image Preparation and Representation}
\subsection{Shared Application Preprocessing}
The deployed preprocessing contract, \texttt{visual-search-rgb-square-v1}, applies
the following operations to gallery and query images:
\begin{enumerate}[leftmargin=1.8em,itemsep=0.3em]
 \item Apply EXIF orientation to the decoded image.
 \item Composite transparency onto a white background and convert to RGB.
 \item Centre the image on a white square whose side is the larger input dimension.
 \item Apply the checkpoint's image processor and run the fixed SigLIP2 image encoder.
 \item Convert the representation to float32 and apply $L_2$ normalisation.
\end{enumerate}
Square padding retains the aspect ratio before the model's own resizing. It
does not segment the gallery. For a SAM2 query, background pixels outside the
accepted mask become white and the result is cropped to the mask's bounding box;
the same square preparation then follows. The implementation produces a
white-background PNG cutout, not a transparent product image.

\begin{figure}[H]
\centering
\begin{tikzpicture}[node distance=5mm]
\node[reportbox,text width=10.9cm] (input) {Original gallery image or accepted query image};
\node[reportbox,text width=10.9cm,below=of input] (rgb) {EXIF correction; white alpha composite; RGB conversion};
\node[reportbox,text width=10.9cm,below=of rgb] (square) {Aspect-preserving square padding, then checkpoint processor};
\node[reportbox,text width=10.9cm,below=of square] (encoder) {Fixed SigLIP2 image encoder; $d=768$};
\node[reportbox,text width=10.9cm,below=of encoder] (norm) {$L_2$ normalisation; float32 vector for inner-product ranking};
\draw[reportarrow] (input)--(rgb);\draw[reportarrow] (rgb)--(square);
\draw[reportarrow] (square)--(encoder);\draw[reportarrow] (encoder)--(norm);
\end{tikzpicture}
\caption{Shared preprocessing and representation contract in the deployed application}
\label{fig:preprocessing}
\end{figure}

\subsection{Index and Artifact Consistency}
The persistent FAISS index is accompanied by a path-mapping JSON file containing
the encoder name, dimension, preprocessing version, metadata schema, image count,
relative paths, and a catalogue checksum. Loading validates these fields against
the installed resources. Incompatible state is rejected so that a later search
or authorised reindex can build a consistent index. The checksum covers image
content and canonical metadata, not only file modification times.

Model weights, index files, query hashes, and protocol settings are recorded
where supported by the relevant experiment. Such records strengthen evidence
about which artifacts were used. They do not retroactively supply provenance
that was absent from an older cached experiment.

\section{Experimental Protocols}
\subsection{Historical Encoder Comparison}
The retained benchmark compares ResNet50, EfficientNet-B3, CLIP, DINOv2, and
SigLIP2. Each uses the same 3,068-image train gallery and 341 validation queries.
Features are normalised and ranked by exact inner product. A retrieved image is
relevant when its category equals the query category.

The saved runs reused gallery and validation embedding caches. The historical
evaluation adapters do not explicitly invoke the deployed square-preprocessing
function. Therefore, these results support the recorded model-selection
comparison, but their zero recorded encoding times cannot be interpreted as
real inference latency, and the values are not presented as a fresh benchmark
of the final application.

\subsection{Historical Full-Image and SAM2 Comparison}
Two retained synthetic protocols use the same construction procedure but
different query partitions. The validation experiment draws ten targets per
category from validation; the final recorded experiment draws ten per category
from test. Each produces 100 two-panel composite images. A target and a
distractor from another class are fitted into separate 512-by-512 white panels
and combined into a 1,024-by-512 image. Seed 42 controls sampling and placement.

The complete composite and the selected-garment cutout are encoded separately
against the same train gallery. A positive point at the centre of the target's
panel is supplied to SAM2. The metric uses the target category. This controlled
construction tests competing garment content, but does not reproduce all the
occlusion, pose, and background complexity of a naturally photographed outfit.
The validation and test scores remain separate because their images differ.

The original scripts reuse existing embeddings and cutouts by existence.
The historical records do not contain all model, code, and original input hashes
from the time of execution. Their stored scores are audited for consistency,
while the report explicitly retains this provenance limitation.

\subsection{Application-Aligned Category Routing}
The category-routing comparison uses the existing original gallery vectors and
newly encodes the clean validation and test queries through
the application's shared query encoder. For each query, unrestricted retrieval and
automatic-centroid filtering use the same representation and gallery. Top-1
correctness is paired by query. No category-confidence threshold is selected or
changed in this comparison.

The original-gallery prediction is the category with highest query-to-centroid
cosine score. A temperature-scaled softmax with $\tau=0.07$ produces the
catalogue-relative display confidence. This value describes the represented
categories and is not a calibrated probability that the first result is correct.

\subsection{Catalogue-Cropping Experiment}
The cropping study creates a separate candidate gallery using a conservative
SAM2-guided bounding crop. Final candidate pixels come from the original image;
the mask is used to locate the crop rather than to whiten every background
pixel. The original-gallery index remains intact. Three configurations are
compared: original embeddings, candidate-crop embeddings, and a weighted
combination of their similarity scores.

Candidate hybrid weights are screened on validation only. The retained mixture
uses 90\% original and 10\% crop similarity. The subsequent comparisons cover
372 clean test images, 100 complete composites, and their 100 selected cutouts.
All configurations within this experiment share its query preparation. These
application-aligned values cannot be pooled with the historical 46\%--75\%
experiment as if they came from one run.

\subsection{Reviewed Alternate-View Identity Pilot}
Twelve visually reviewed groups were available from earlier catalogue review.
Five pairs were conservatively excluded before inference because independence
of the views was uncertain. Seven groups remained. A fixed rule selected one
view per group as query and the other as the accepted target. All seven queries,
and any byte-identical or decoded-pixel-identical copies, were removed from the
gallery before any condition. The retained gallery contained 3,061 images.
Category centroids were recomputed from that retained gallery only.

The forward pilot compares automatic and unrestricted search. A separately
specified reverse pilot swaps the two views and compares full images with SAM2
cutouts on the same reverse gallery. The reverse prompts were selected visually
by the assistant and sealed before reverse retrieval results, without subsequent
point refinement. They were not user-approved reference masks. A SAM2 failure
would remain a failure in the denominator; none occurred in the seven cases.

These are selected development cases, not independent consumer photographs or
seller-verified SKU labels. The reverse direction reuses the same seven groups
and is not counted as seven additional independent products. The protocol
prevents direct self-match but cannot establish dataset-wide identity accuracy.

\section{Evaluation Measures}
Let $Q$ be the query set, $R_q$ the number of gallery images judged relevant to
query $q$, and $r_q(j)\in\{0,1\}$ the relevance of the result at rank $j$.
Relevance is category equality in the main experiments and reviewed-group
membership in the identity pilot. Following the ranked-retrieval concepts in
\cite{Manning2008IR}, the implemented cutoff conventions are:
\begin{align}
 P_q@K &= \frac{1}{K}\sum_{j=1}^{K} r_q(j), \\
 \operatorname{Hit}_q@K &= \mathbb{1}\!\left[\sum_{j=1}^{K}r_q(j)>0\right], \\
 \operatorname{AP}_q@K &= \frac{\sum_{j=1}^{K}(P_q@j)r_q(j)}{\min(R_q,K)}, \\
 \operatorname{mAP}@K &= \frac{1}{|Q|}\sum_{q\in Q}\operatorname{AP}_q@K, \\
 \operatorname{MRR}@K &= \frac{1}{|Q|}\sum_{q\in Q}\frac{\mathbb{1}[\rho_q\leq K]}{\rho_q}.
\end{align}
Here $\rho_q$ is the first relevant rank; its contribution is zero if no relevant
result appears within the cutoff. AP is defined as zero for an empty relevant
set. Mean P@K averages the per-query values. At $K=1$, category P@1 and Hit@1
coincide. Identity Top-$K$ in the pilot is a hit rate over reviewed target groups,
not precision among all $K$ returned images.

For paired Top-1 comparisons, let $b$ count queries corrected by the candidate
and $c$ count queries made incorrect. The observed gain is
\begin{equation}
 \Delta_{\mathrm{Top1}}=\frac{b-c}{|Q|}.
\end{equation}
The saved comparison scripts resample paired query differences 10,000 times
with seed 42 to obtain percentile 95\% bootstrap intervals. They also report
two-sided exact McNemar tests over the discordant pairs. These describe sampling
uncertainty under the recorded query set; they do not repair class-label or
selection bias. Statistical results are interpreted alongside the effect size
and experimental design.

\section{Software Development and Verification Method}
A baseline, an isolated candidate change, comparison, and justified promotion
structure development. The original index is retained when evidence does not
support a replacement. Unit tests cover validation, metadata, routing, revisions,
authentication, and backups. JavaScript tests use controlled responses to exercise
selection and request state; HTTP checks use actual inference, and browser
observations examine the visible workflow.

Final-report verification reruns these checks using an independent application
copy with matching active code and separate account/product databases. Records
state their date, command, scope, and limitations. A functional test is not a
retrieval query and contributes no evidence of general model accuracy.

\section{Ethical and Reproducibility Considerations}
Dataset and model use retains upstream notices. Research fixtures are
identified separately from user uploads. Monitoring stores query hashes and
search metadata, not raw search photographs; optional user notes still require
normal local data care. Results retain their denominator, relevance definition,
gallery, and preprocessing. Re-reading a historical result is not described as
new inference, and an exact-product claim requires appropriate identity evidence.

\section{Summary}
The method separates category experiments, reviewed identity pilots, and software
verification. These protocols constrain the interpretation of results and lead
to the operating requirements analysed in the next chapter.
